\documentclass[10pt,a4paper]{amsart}
\usepackage[margin=19mm]{geometry}
\usepackage{amsmath,amssymb,mathtools}
\usepackage[scr=rsfs,cal=euler]{mathalfa}
\usepackage{xfrac}
\usepackage[unicode,hidelinks,bookmarks=false]{hyperref}
\newcommand{\ie}{i.e.\ }
\newcommand{\cf}{cf.\ }
\newcommand{\resp}{respectively\ }
\newcommand{\Subsection}{\subsection}
\providecommand{\nobreakdash}{\nobreak\hbox{-}}
\setlength{\emergencystretch}{2em}
\multlinegap0pt
\usepackage{bm}
\usepackage{bbm}
\usepackage{dsfont}
\usepackage{xspace}
\usepackage{cancel}
\usepackage{mathtools}
\usepackage{amsthm}
\makeatletter
\@ifundefined{theorem}{\newtheorem{theorem}{Theorem}}{}
\makeatother
\title[Sparse Polynomial Regression under Anomalous Data]{Sparse Polynomial Regression under Anomalous Data}
\author{Roozbeh Abolpour, Mohammad Reza Hesamzadeh, Maryam Dehghani}
\date{Source: arXiv:2508.18199v1 (CC BY 4.0)}
\begin{document}
\maketitle

\noindent\emph{Source and licence.} Sparse Polynomial Regression under Anomalous Data. Version arXiv:2508.18199v1 (CC BY 4.0), distributed under the Creative Commons Attribution 4.0 International licence, as recorded in the arXiv licence metadata for this exact version and shown on the abstract page https://arxiv.org/abs/2508.18199v1. Full rights notice: licenses/arxiv-2508.18199-CC-BY-4.0.txt.\par\smallskip

\noindent\textit{Editorial scope.} Counted unit: the Theorem whose source label is \texttt{thm4\_2} in the article (source0.tex lines 318--320, source label \texttt{thm4\_2}), delivered together with the article's own complete original proof of it (source0.tex lines 321--376, one \texttt{proof} environment of 56 lines). No mathematics is written, repaired or re-derived: statement and proof are the article's own text line for line. Required context retained uncounted: eq8 (lines 162--175); eq26\_d (lines 244--267); eq26\_e (lines 244--267); eq26\_h (lines 244--267); eq26\_j (lines 244--267); eq26\_k (lines 261--265). Every definition, display and label of those lines is retained unchanged. Omitted from this package and not redistributed: 1--161, 176--243, 266--317, 377--862. Nothing inside a retained excerpt is omitted or altered: every retained body is delivered exactly as the article's lines are written, so no omission receipt of any kind accompanies this package. Citations: the retained excerpts carry no citation command at all, so no bibliography entry of the article is transcribed and no reference is redistributed. Reference adaptation: none was needed and none was made; every reference inside the delivered unit resolves inside it, so no unresolved reference or citation is produced. Printed slips kept literal: no printed word, symbol, hypothesis or number of the counted statement or of its proof was altered. Layout and class: the article's own class is \texttt{svjour3} with packages including natbib, graphicx, multirow, amsmath, amssymb, amsfonts, mathrsfs, appendix, xcolor, textcomp, manyfoot, booktabs, algorithm, algpseudocode; the wrapper is the lane's pinned \texttt{amsart} editorial wrapper at 10pt on A4, which restates the theorem-like environments the retained text needs and adds the article's own macro definitions that the retained text needs (none), with extra wrapper packages (bm, bbm, dsfont, xspace, cancel, mathtools). The delivered numbering is the unit's own. Assets: no retained span contains a figure environment, a graphics-inclusion command, a PDF-page inclusion, a table or a TikZ picture; no image file is copied into this package, the package declares no asset dependency and no image file is redistributed with it. Licence: version arXiv:2508.18199v1 of this article is distributed under the Creative Commons Attribution 4.0 International licence, as recorded in the arXiv licence metadata for this exact version and shown on the abstract page https://arxiv.org/abs/2508.18199v1. A full rights notice is delivered at licenses/arxiv-2508.18199-CC-BY-4.0.txt. Characters: every retained excerpt is pure ASCII, so the delivered engine, XeLaTeX with its default Latin Modern text font, reports no missing character.
\begin{subequations} \label{eq8}
	\begin{align}
        \underset{\Omega_4}{\min}&\quad \frac{\hat{\gamma}}{\hat{v}} \label{eq8_a} \\ \displaybreak[0]
		\text{s.t.} &\quad -M\hat{v}+\frac{\sqrt{\rho-1}}{\sqrt{\rho}}M\hat{b}_k+\left|y^{(k)}\hat{v}-\sum_{\alpha \in \Gamma_d} \hat{c}_\alpha {x^{(k)}}^{\alpha} \right| \le \hat{\gamma},\quad \forall k \in \{1,\dots,N\}, \label{eq8_b} \\ \displaybreak[0]
        &\quad -\frac{\sqrt{\rho-1}}{\sqrt{\rho}} M \hat{s}_\alpha \le \hat{c}_\alpha \le \frac{\sqrt{\rho-1}}{\sqrt{\rho}} M \hat{s}_\alpha,\quad \forall \alpha \in \Gamma_d, \label{eq8_c} \\ \displaybreak[0]
        &\quad 0 \le \hat{s}_\alpha \le \frac{\sqrt{\rho}}{\sqrt{\rho - 1}} \hat{v},\quad \forall \alpha \in \Gamma_d, \label{eq8_d} \\ \displaybreak[0]
        &\quad 0 \le \hat{b}_k \le \frac{\sqrt{\rho}}{\sqrt{\rho - 1}} \hat{v},\quad\forall k \in \{1,\dots,N\}, \label{eq8_e} \\ \displaybreak[0]
        &\quad \sum_{k=1}^N \hat{b}_k = l_b  \frac{\sqrt{\rho}}{\sqrt{\rho - 1}}, \label{eq8_f} \\ \displaybreak[0]
        &\quad \sum_{\alpha \in \Gamma_d} \hat{s}_\alpha = l_m \frac{\sqrt{\rho}}{\sqrt{\rho - 1}}, \label{eq8_g} \\ \displaybreak[0]
        &\quad \hat{v} = \frac{\sqrt{\rho - 1}}{\sqrt{\rho(m_d + N + 1) - 1}}, \label{eq8_h} \\ \displaybreak[0]
        &\quad \frac{(\rho - 1)(\|\hat{s}\|^2 + \|\hat{b}\|^2)}{\rho} + \frac{\hat{v}}{\sqrt{\rho} \sqrt{\rho - 1}} \left( \sum_{\alpha \in \Gamma_d} \hat{s}_\alpha + \sum_{k=1}^N \hat{b}_k \right) + \hat{v}^2 \le 1, \label{eq8_i} \\ \displaybreak[0]
        &\quad \|\hat{s}\|^2 + \|\hat{b}\|^2 + \hat{v}^2 \ge 1, \label{eq8_j}
	\end{align}
\end{subequations}

	\begin{align}
        \underset{\Omega_5}{\text{min}}&\quad~\frac{\hat{\gamma}}{\hat{v}}, \label{eq26_a} \\ \displaybreak[0]
		\text{s.t.}&\quad -M\hat{v}+\frac{\sqrt{\rho-1}}{\sqrt{\rho}}M\hat{b}_k+\left|y^{(k)}\hat{v}-\sum_{\alpha \in \Gamma_d} \hat{c}_\alpha {x^{(k)}}^{\alpha} \right| \le \hat{\gamma},\quad \forall k \in \{1,\dots,N\}, \label{eq26_b} \\ \displaybreak[0]
        &\quad -\frac{\sqrt{\rho-1}}{\sqrt{\rho}} M \hat{s}_\alpha \le \hat{c}_\alpha \le \frac{\sqrt{\rho-1}}{\sqrt{\rho}} M \hat{s}_\alpha,\quad \forall \alpha \in \Gamma_d, \label{eq26_c} \\ \displaybreak[0]
        &\quad 0 \le \hat{s}_\alpha \le \frac{\sqrt{\rho}}{\sqrt{\rho - 1}} \hat{v},\quad \forall \alpha \in \Gamma_d, \label{eq26_d} \\ \displaybreak[0]
        &\quad 0 \le \hat{b}_k \le \frac{\sqrt{\rho}}{\sqrt{\rho - 1}} \hat{v},\quad \forall k \in \{1,\dots,N\}, \label{eq26_e} \\ \displaybreak[0]
        &\quad \sum_{k=1}^N \hat{b}_k = l_b \frac{\sqrt{\rho}}{\sqrt{\rho - 1}}, \label{eq26_f} \\ \displaybreak[0]
        &\quad \sum_{\alpha \in \Gamma_d} \hat{s}_\alpha = l_m  \frac{\sqrt{\rho}}{\sqrt{\rho - 1}}, \label{eq26_g} \\ \displaybreak[0]
        &\quad \hat{v} = \frac{\sqrt{\rho - 1}}{\sqrt{\rho(m_d + N + 1) - 1}}, \label{eq26_h} \\ \displaybreak[0]
        &\quad \frac{(\rho - 1)(\|\hat{s}\|^2 + \|\hat{b}\|^2)}{\rho} + \frac{\hat{v}}{\sqrt{\rho} \sqrt{\rho - 1}} \left( \sum_{\alpha \in \Gamma_d} \hat{s}_\alpha + \sum_{k=1}^N \hat{b}_k \right) + \hat{v}^2 \le 1, \label{eq26_i} \\ \displaybreak[0]
        &\quad \operatorname{trace}(\hat{\chi}) \ge 1, \label{eq26_j} \\ \displaybreak[0]
        &\quad \hat{\chi} \succeq 
        \begin{bmatrix}
            \hat{s} \\
            \hat{b} \\
            \hat{v}
        \end{bmatrix}
        \begin{bmatrix}
            \hat{s} \\
            \hat{b} \\
            \hat{v}
        \end{bmatrix}^T, \label{eq26_k} \\ \displaybreak[0]
        &\quad \operatorname{rank}(\hat{\chi}) = 1, \label{eq26_l}
	\end{align}

\begin{theorem} \label{thm4_2}
The linear-based relaxation of FP \eqref{eq8} is tighter than the SDC-based relaxation of FP \eqref{eq8}.
\end{theorem}

\begin{proof}
Let $\left(\hat{c}_l, \hat{s}_l, \hat{b}_l, \hat{\gamma}_l, \hat{v}_l\right)$ be an arbitrary feasible solution to the linear-based relaxation of FP \eqref{eq8}. Then, by definition, the following inequality must hold
\begin{equation} \label{eq:sum1}
\sum_{\alpha \in \Gamma_d} \hat{s}_{l,\alpha} + \sum_{k=1}^N \hat{b}_{l,k} + \hat{v}_l \ge 1.
\end{equation}

We now prove that this point also belongs to the feasible region of the SDC-based relaxation by constructing a suitable matrix $\hat{\chi}_l$ and verifying the SDC constraints.

\textbf{Case 1:} Assume $\|\hat{s}_l\|^2 + \|\hat{b}_l\|^2 + \hat{v}_l^2 \ge 1$. Let
\begin{equation}
\hat{\chi}_l = 
\begin{bmatrix}
\hat{s}_l \\
\hat{b}_l \\
\hat{v}_l
\end{bmatrix}
\begin{bmatrix}
\hat{s}_l \\
\hat{b}_l \\
\hat{v}_l
\end{bmatrix}^\top.
\end{equation}

Then, $\hat{\chi}_l \succeq \begin{bmatrix} \hat{s}_l \\ \hat{b}_l \\ \hat{v}_l \end{bmatrix} \begin{bmatrix} \hat{s}_l \\ \hat{b}_l \\ \hat{v}_l \end{bmatrix}^\top$ holds trivially and $\operatorname{trace}(\hat{\chi}_l)=\|\hat{s}_l\|^2 + \|\hat{b}_l\|^2 + \hat{v}_l^2 \ge 1$ by the case's assumption. Thus, tuple $(\hat{c}_l, \hat{s}_l, \hat{b}_l, \hat{\gamma}_l, \hat{v}_l, \hat{\chi}_l)$ satisfies the SDC constraints \eqref{eq26_k} and \eqref{eq26_j}, and is a feasible point of the SDC-based relaxation.

\textbf{Case 2:} Assume $\|\hat{s}_l\|^2 + \|\hat{b}_l\|^2 + \hat{v}_l^2 < 1$. Define
\begin{equation} \label{eq143}
z = \begin{bmatrix}
\hat{s}_l \\
\hat{b}_l \\
\hat{v}_l
\end{bmatrix}, \quad
e = \mathbf{1} \in \mathbb{R}^{m_d + N + 1}, \quad
\hat{\chi}_l = \frac{1}{2}(z e^\top + e z^\top).
\end{equation}

Due to (\ref{eq26_d}), (\ref{eq26_e}), and(\ref{eq26_h}), all entries in $z$ are non-negative. Using this fact and inequality $\|\hat{s}_l\|^2 + \|\hat{b}_l\|^2 + \hat{v}_l^2 < 1$, the following relation can be directly obtained.
\begin{equation} \label{eq144}
\sum_{i=1}^{m_d+N+1} z_i^2 \le \sum_{i=1}^{m_d+N+1} z_i \quad \Rightarrow \quad \|z\|^2 \le \sum_{i=1}^{m_d+N+1} z_i.    
\end{equation}

Then $\hat{\chi}_l$ is symmetric by construction. We now show that $\hat{\chi}_l \succeq zz^\top$ by verifying that
\begin{equation}
y^\top \hat{\chi}_l y \ge y^\top zz^\top y, \quad \forall y \in \mathbb{R}^{m_d + N + 1}.
\end{equation}

(a) If $y \perp z$, then $y^T z = 0$, and the inequality
$y^\top \hat{\chi}_l y = \frac{1}{2}(y^\top z)(e^\top y) + \frac{1}{2}(y^\top e)(z^\top y) = 0$  holds with equality.

(b) If $y = z$, then $y^\top \hat{\chi}_l y = \frac{1}{2}(z^\top z)(e^\top z) + \frac{1}{2}(z^\top e)(z^\top z) = (z^\top z)(e^\top z) = \|z\|^2 \cdot \sum_i z_i\ge \|z\|^4=y^T(zz^T)y$ based on (\ref{eq144}). 

Together, these cases show that $\hat{\chi}_l \succeq zz^\top$, i.e., SDC constraint \eqref{eq26_k} holds. Finally, by the linear constraint \eqref{eq:sum1}, we have $\operatorname{trace}(\hat{\chi}_l) = z^\top e = \sum_{\alpha \in \Gamma_d} \hat{s}_{l,\alpha} + \sum_{k=1}^N \hat{b}_{l,k} + \hat{v}_l \ge 1
$. Therefore, constraint \eqref{eq26_j} is also satisfied. Thus, the tuple $\left(\hat{c}_l, \hat{s}_l, \hat{b}_l, \hat{\gamma}_l, \hat{v}_l, \hat{\chi}_l\right)$ is feasible for the SDC-based relaxation. 

Since this holds for any feasible point of the linear-based relaxation, it is tighter than the SDC-based relaxation. \qed
\end{proof}

\end{document}
